\scanpage{71}
\uline{Théorème de finitude sur Picard} (Commentaires

Kleimann prouve, pour $X/k$ \uline{projectif}, les critères \textcircled{\scriptsize viii} de $\tau$-équivalence à $0$, sans utiliser les fourbis délicats (i) (ii) \ldots. Pour récupérer le cas $X/k$ propre, il doit utiliser (i) (via lemme de Chow). Avec le même grain de sel, il prouve \textcircled{\scriptsize v} que $\underline{\mathrm{Pic}}^\tau_{X/S}$ est qu. cpct. sur $S$. [\uline{En fait, il prouve directement que $\underline{\mathrm{Pic}}^{\mathrm{num}}_{X/S}$ est quasi-compact\ldots.}].

Utilisant le th. de l'index pour une surface projective non singulière, Kleimann prouve de plus le ``critère d'équivalence'' sous la forme suivante : Si $X$ est projective $/k$, ttes \unsure{ses|les|es} composantes de dim $\geq3$, et $Y$ section hyperplane, d'où $f^*\colon\underline{\mathrm{Pic}}_{\deleted{\unsure{Y}}X/k}\unsure{-|\to}\underline{\mathrm{Pic}}_{Y/k}$, alors $f^{*^{-1}}(\underline{\mathrm{Pic}}^\tau_{Y/k})=\underline{\mathrm{Pic}}^\tau_{X/k}$, i.e. si $k$ \unsure{sp.|sep.|alg.} clos, $\mathrm{NSL}(X)\to\mathrm{NSL}(Y)$ est injectif. Ce\deleted{\unsure{ci}}\inserted{ qui précède}, et l'``inégalité'' de \uline{Picard. Igusa}, \inserted{et le th. de \uline{Néron-Sévéri} \deleted{\unsure{pour}} pour les surfaces proj. non sing., et \unsure{résolution} des surfaces \inserted{de \uline{Abhyankar}}} \unsure{suffisent|suffit} à entraîner \textcircled{\scriptsize vi} que si $X/S$ est propre de prés. finie sur \deleted{$k$} $S$ quasi-compact, les groupes de N.S. des fibres géométriques sont de \deleted{rang} \inserted{type fini} ``majoré'' \deleted{\unsure{et de}} (i.e. groupe de torsion d'ordre majoré, et rang majoré). [Dans le cas $X/S$ projectif, cela sort donc des dévissages de Kleimann, sans utiliser mes fourbis]. On récupère aussi ainsi \textcircled{\scriptsize vii}.

\uline{TSVP}
% archive sheet number: 70

